\section*{Chapter \ref{ch:m1:options-market-structure} --- Options Market Structure}

\begin{solution}{exo:m1:options-market-structure:1}
C1 receives 10 and C2 receives 2: customers in time order, and the order is
exhausted before any professional is reached.
\end{solution}

\begin{solution}{exo:m1:options-market-structure:2}
Customers: 10 and 5, leaving 100. The designated maker's pro-rata share would
be $100 \times 100/500 = 20$; its entitlement is 40. The remaining 60 are
shared over the 400 lots of MM1, F and MM2: $\lfloor 60 \times 100/400\rfloor
= 15$, $\lfloor 60 \times 50/400\rfloor = 7$, $\lfloor 60 \times 250/400\rfloor
= 37$; the \omterm{def:m1:us-equity-market-structure:lots}{odd lot} goes to the oldest, MM1: 16, 7, 37.
\end{solution}

\begin{solution}{exo:m1:options-market-structure:3}
Valid: 2.43, 2.99 (cents below \$3.00), 3.05 and 12.10 (five-cent steps from
\$3.00). Invalid: 3.02 and 12.12. Around 3.10 the minimum spread is 0.05, that
is 1.6\% of the price.
\end{solution}

\begin{solution}{exo:m1:options-market-structure:4}
$311 \times 10^9 / 23\,400\,\text{s} = 13.3$ million messages a second on
average. Peak: 13.575 million per 100 milliseconds, a rate of 135.8 million a
second, 10.2 times the average. Bandwidth: 4.403 gigabits per 100 milliseconds
is 44 gigabits a second for one stream; two streams with 10\% for
retransmissions: 97 gigabits a second.
\end{solution}

\begin{solution}{exo:m1:options-market-structure:5}
Nobody: the initiator fills all 100. One responder at the stop: 50. Three at
the stop: 40, the responders sharing 60. One at 2.49 for 60 and two at 2.50:
the 60 improved lots go to the improver; at the stop the initiator may take
up to 40\% of the original order, which is the whole remaining 40.
\end{solution}

\begin{solution}{exo:m1:options-market-structure:6}
Leg by leg: pay 2.50, receive 0.90: 1.60. Midpoint: $2.45 - 0.95 = 1.50$. The
package's displayed market is 1.40 at 1.60, 0.20 wide. The spread's delta and
vega are a fraction of either leg's, so a \omterm{def:m1:what-a-trading-firm-does:market-maker}{market maker}'s risk in quoting it
is small; and the errors in its two leg valuations largely cancel. A complex
book lets that be expressed as one price, a cent or two from 1.50.
\end{solution}

\begin{solution}{exo:m1:options-market-structure:7}
One responder on average: the initiator keeps 47.8\% and 39.7\% of the order
is improved. Three: 12.7\% and 77.5\%. The third responder almost doubles the
share of the customer's order that receives a better price, and takes three
quarters of what the initiator was keeping.
\end{solution}

\begin{solution}{exo:m1:options-market-structure:8}
Time priority does not exist among professionals here. Any customers at the
price are filled first; then the 100 lots are shared by size: the 5-lot bid's
share is $\lfloor 100 \times 5/455\rfloor = 1$ lot, plus possibly the \omterm{def:m1:us-equity-market-structure:lots}{odd lot}
left by rounding: 2 lots at best, the \omterm{def:m1:what-a-trading-firm-does:market-maker}{market makers} 98. To be filled in such
a market a professional must show size, improve the price, or respond to
auctions; arriving first earns nothing.
\end{solution}

\begin{solution}{pb:m1:options-market-structure:1}
\textbf{1.} 100, 50, 40, 40 (60 improved lots go elsewhere and the initiator
takes the remaining 40, its 40\% of the original) and 0.
\textbf{2.} $10 + 12.5 + 14 + 8 + 0 = 44.5$ lots.
\textbf{3.} It sells at 2.50 what it values at 2.44: 6 cents a share, \$6 a
contract.
\textbf{4.} $44.5 \times 6 = \$267$.
\textbf{5.} One cent on 60 lots in 20\% of cases and on 100 lots in 10\%:
$0.2 \times 60 + 0.1 \times 100 = 22$ lots improved by \$1 each, \$0.22 per
contract on average over the order.
\textbf{6.} 2.50, the displayed offer, if 100 lots are there. The \omterm{def:m1:retail-flow-and-wholesaling:wholesaler}{wholesaler}
would earn nothing on the trade itself.
\textbf{7.} $267 - 100 \times 0.30 = \$237$.
\textbf{8.} The customer pays at most 2.48: two cents of improvement for
certain. The \omterm{def:m1:retail-flow-and-wholesaling:wholesaler}{wholesaler}'s edge falls to 4 cents a share. Responders now need
to bid 2.47 to improve, which few will at a value of 2.44 to 2.45: the
\omterm{def:m1:retail-flow-and-wholesaling:wholesaler}{wholesaler} keeps a larger share of a thinner margin. The stop price is the
\omterm{def:m1:retail-flow-and-wholesaling:wholesaler}{wholesaler}'s choice and its main instrument.
\textbf{9.} At 2.50 and 2.49 (edges of 5 and 4 cents); at 2.46 it has a cent.
In time means receiving the auction message, recognising the series,
repricing it from the current share price and volatility, checking risk and
answering, all well inside 100 milliseconds less two network trips: an
automated, co-located pricing engine.
\textbf{10.} Not as a resting quote competing with responses on equal terms:
priority customers on the book are honoured, professional quotes are not
guaranteed a share of an order that was brought to the auction already
paired. To take part it must send a response.
\textbf{11.} Responding reveals its price for that series to the initiator
and to every participant who sees the auction, a hundred times a second; and
the orders brought to auctions are selected. It may still respond, because
auctioned flow is the best flow there is.
\textbf{12.} Orders from professionals, large institutional orders worked by
algorithms, and orders of firms reacting to news: the flow that reaches
displayed quotes unaccompanied is, on average, better informed.
\textbf{13.} Wider. A displayed quote is an option given to the fastest and
best-informed participants; it should be priced for them.
\textbf{14.} It is the segmentation of \cref{ch:m1:retail-flow-and-wholesaling}
carried out on an exchange: uninformed flow is separated, priced tightly and
shared among those with a claim on it; the lit quote is left to serve
everyone else.
\textbf{15.} For: several well-capitalised firms see every auction and
respond by machine; a tenth of a second is long for them; improvement does
occur. Against: only co-located firms can respond; the initiator chooses the
stop and keeps up to half whatever happens; a responder that merely matches
gets a minority share, so the incentive to improve is weak unless it expects
to win the whole order.
\textbf{16.} Whoever trades against the \omterm{def:m1:what-a-trading-firm-does:market-maker}{market makers}' quotes. The fee is a
cost of market making in customer flow, recovered in the spread; customers
receive it back, in part, through their brokers' commission-free pricing.
\textbf{17.} To serve both constituencies at once, makers on one exchange
and takers on another, to obtain more entitlements and more market-data and
connectivity revenue, and because each exchange licence is a separate place
in every broker's routing table.
\textbf{18.} Less to improve: with a spread of one or two cents the stop is
already near fair value, auctions would improve by less, and \omterm{def:m1:retail-flow-and-wholesaling:wholesaler}{wholesalers}'
margins would fall, as would displayed size.
\textbf{19.} \textbf{44.5 lots of the 100.}
\textbf{20.} A few dollars a contract to the firm that fills it, which is why
brokers are paid for it, exchanges build mechanisms around it, and the
displayed market is what remains.
\end{solution}

\begin{solution}{iq:m1:options-market-structure:1}
One clearing house and many exchanges, as in equities, but: hundreds of
times more instruments, almost all liquidity from \omterm{def:m1:what-a-trading-firm-does:market-maker}{market makers}' quotes
rather than natural orders; \omterm{def:m1:options-market-structure:customer}{customer priority} and \omterm{def:m1:matching-algorithms-and-implied-spreads:prorata}{pro-rata allocation} with
market-maker entitlements on many exchanges; fees that depend on the
trader's capacity; \omterm{def:m1:options-market-structure:auction}{price-improvement auctions} as the main channel for retail
flow instead of off-exchange \omterm{def:m1:retail-flow-and-wholesaling:wholesaler}{internalisation}; \omterm{def:m1:options-market-structure:complex}{complex order books}; and a
consolidated feed of hundreds of billions of messages a day.
\iqlookfor{quote-driven, capacity-dependent rules, auctions.}
\end{solution}

\begin{solution}{iq:m1:options-market-structure:2}
Public customers' resting orders at the best price are filled before any
professional's, regardless of time. Exchanges grant it to attract customer
orders, which professionals value because they are uninformed; the customer
flow then attracts the \omterm{def:m1:what-a-trading-firm-does:market-maker}{market makers}, whose quotes make the exchange's
market.
\iqlookfor{the two-sided-platform logic.}
\end{solution}

\begin{solution}{iq:m1:options-market-structure:3}
Initiator: choose which customer orders to auction and at what stop; the
lower the stop on a buy, the better for the customer and the thinner the
edge, but the larger the share kept. Submit the pair; the exchange guarantees
the customer the stop. Outcome: better-priced responses first; at the stop,
customers on the book, then up to 50\% or 40\% to the initiator, then the
rest. Responder: receive the message, price the series now, decide between
matching (a pro-rata share of what the initiator leaves) and improving (first
claim), and answer within the period.
\iqlookfor{the stop price as the strategic variable; match versus improve.}
\end{solution}

\begin{solution}{iq:m1:options-market-structure:4}
Volume and burstiness: hundreds of billions of messages a day with peaks ten
times the average within 100 milliseconds, tens of gigabits a second per
stream, two streams to arbitrate, gap detection and retransmission. Filtering
as early as possible (most series are irrelevant to any one strategy),
ideally in hardware or kernel-bypass receivers; symbol mapping for over a
million series that change daily; building books per series per exchange;
and conflation policies for slow consumers that never reorder or invent
prices.
\iqlookfor{peak-to-average design, early filtering, line arbitration.}
\end{solution}

\begin{solution}{iq:m1:options-market-structure:5}
All of a maker's quotes in a class are one opinion about the underlying; if
that opinion is stale, an informed trader can hit hundreds of them before
the maker reprices. Protection caps the damage by cancelling everything after
a number of fills in a window. The cost: displayed size is partly illusory,
since it disappears after the first few fills, so a taker of size sees
quotes fade; and the mechanisms can themselves withdraw liquidity abruptly
in fast markets.
\iqlookfor{correlation of quotes as the reason; fading as the cost.}
\end{solution}

\begin{solution}{iq:m1:options-market-structure:6}
The least informed flow is diverted into auctions and allocated by
entitlement, so what reaches my displayed quote unaccompanied is adversely
selected: professionals, algorithms, news. I should quote the book wider
than my fair spread for average flow, size it for toxic flow, protect it
with fast underlying-driven repricing and risk limits, and compete for the
benign flow where it actually goes: respond to auctions, obtain a
designated-maker role where the entitlement pays for the obligations, or
negotiate directed flow. Measure each channel's markouts separately.
\iqlookfor{\omterm{def:m1:what-a-trading-firm-does:adverse-selection}{adverse selection} of the residual lit flow; going to where the
flow is.}
\end{solution}
